\documentclass[a4paper]{article}
\usepackage[utf8]{inputenc}
\usepackage[T1]{fontenc}
\usepackage[spanish]{babel}
\usepackage{a4wide}
\usepackage{amsmath}
\usepackage{amssymb,amsfonts,textcomp}
\usepackage{color}
\usepackage{array}
\usepackage{supertabular}
\usepackage{hhline}
\usepackage{hyperref}
\usepackage{fancyhdr}
\usepackage[table]{xcolor}
%\hypersetup{pdftex, colorlinks=true, linkcolor=blue, citecolor=blue, filecolor=blue, urlcolor=blue, pdftitle=, pdfauthor=, pdfsubject=, pdfkeywords=}
\usepackage[pdftex]{graphicx}
% footnotes configuration
\makeatletter
\renewcommand\thefootnote{\arabic{footnote}}
\makeatother
% Text styles
\newcommand\textstyleStrongEmphasis[1]{\textbf{#1}}
% Outline numbering
\setcounter{secnumdepth}{0}
\makeatletter
\newcommand\arraybslash{\let\\\@arraycr}
\makeatother
% List styles
\newcommand\liststyleWWNumxxix{%
\renewcommand\labelitemi{{}-}
\renewcommand\labelitemii{o}
\renewcommand\labelitemiii{[F0A7?]}
\renewcommand\labelitemiv{[F0B7?]}
}
\newcommand\liststyleLi{%
\renewcommand\labelitemi{{\textbullet}}
\renewcommand\labelitemii{${\circ}$}%
\renewcommand\labelitemiii{${\blacksquare}$}
\renewcommand\labelitemiv{{\textbullet}}
}
\newcommand\liststyleLii{%
\renewcommand\labelitemi{{\textbullet}}
\renewcommand\labelitemii{${\circ}$}
\renewcommand\labelitemiii{${\blacksquare}$}
\renewcommand\labelitemiv{{\textbullet}}
}
\newcommand\liststyleLiii{%
\renewcommand\labelitemi{{\textbullet}}
\renewcommand\labelitemii{${\circ}$}
\renewcommand\labelitemiii{${\blacksquare}$}
\renewcommand\labelitemiv{{\textbullet}}
}
\newcommand\liststyleLiv{%
\renewcommand\labelitemi{{\textbullet}}
\renewcommand\labelitemii{${\circ}$}
\renewcommand\labelitemiii{${\blacksquare}$}
\renewcommand\labelitemiv{{\textbullet}}
}
\newcommand\liststyleLv{%
\renewcommand\labelitemi{{\textbullet}}
\renewcommand\labelitemii{${\circ}$}
\renewcommand\labelitemiii{${\blacksquare}$}
\renewcommand\labelitemiv{{\textbullet}}
}
\newcommand\liststyleLvi{%
\renewcommand\labelitemi{{\textbullet}}
\renewcommand\labelitemii{${\circ}$}
\renewcommand\labelitemiii{${\blacksquare}$}
\renewcommand\labelitemiv{{\textbullet}}
}
\newcommand\liststyleLvii{%
\renewcommand\labelitemi{{\textbullet}}
\renewcommand\labelitemii{{\textbullet}}
\renewcommand\labelitemiii{{\textbullet}}
\renewcommand\labelitemiv{{\textbullet}}
}
\newcommand\liststyleLviii{%
\renewcommand\labelitemi{{\textbullet}}
\renewcommand\labelitemii{{\textbullet}}
\renewcommand\labelitemiii{{\textbullet}}
\renewcommand\labelitemiv{{\textbullet}}
}
\newcommand\liststyleLix{%
\renewcommand\labelitemi{{\textbullet}}
\renewcommand\labelitemii{{\textbullet}}
\renewcommand\labelitemiii{{\textbullet}}
\renewcommand\labelitemiv{{\textbullet}}
}
\newcommand\liststyleLx{%
\renewcommand\labelitemi{{\textbullet}}
\renewcommand\labelitemii{{\textbullet}}
\renewcommand\labelitemiii{{\textbullet}}
\renewcommand\labelitemiv{{\textbullet}}
}
% Page layout (geometry)
\setlength\voffset{-1in}
\setlength\hoffset{-1in}
\setlength\topmargin{2cm}
\setlength\oddsidemargin{2cm}
\setlength\textheight{24.203001cm}
\setlength\textwidth{17.001cm}
\setlength\footskip{0.0cm}
\setlength\headheight{0.998cm}
\setlength\headsep{0.499cm}
% Footnote rule
\setlength{\skip\footins}{0.119cm}
\renewcommand\footnoterule{\vspace*{-0.018cm}\setlength\leftskip{0pt}\setlength\rightskip{0pt plus 1fil}\noindent\textcolor{black}{\rule{0.25\columnwidth}{0.018cm}}\vspace*{0.101cm}}
% Pages styles
\makeatletter
\newcommand\ps@Standard{
  \renewcommand\@oddhead{[Warning: Draw object ignored]}
  \renewcommand\@evenhead{\@oddhead}
  \renewcommand\@oddfoot{}
  \renewcommand\@evenfoot{}
  \renewcommand\thepage{\arabic{page}}
}
\newcommand\ps@FirstPage{
  \renewcommand\@oddhead{}
  \renewcommand\@evenhead{}
  \renewcommand\@oddfoot{}
  \renewcommand\@evenfoot{}
  \renewcommand\thepage{\arabic{page}}
}
\makeatother

\setlength\tabcolsep{1mm}
\renewcommand\arraystretch{1.3}
\title{}
\author{}
\date{2013-06-03}

\renewcommand{\baselinestretch}{1.2}

\begin{document}

\pagestyle{empty}
\setcounter{page}{0}

\includegraphics[width=3cm]{logo-uclm.jpg}
\hfill
\includegraphics[width=3cm]{logo-esii.png}

{\centering \par}
\begin{center}
\includegraphics[width=2.96cm,height=3.522cm]{logo-conciti.png}
\end{center}

\vskip 3em
{\centering\bfseries\large
UNIVERSIDAD DE CASTILLA-LA MANCHA
\par}

\vskip 3em
{\centering\bfseries\large
ESCUELA SUPERIOR DE INFORMÁTICA
\par}

{\centering\bfseries\large
Departamento de Sistemas Informáticos\footnote{DEPARTAMENTO DE SISTEMAS INFORMÁTICOS o DEPARTAMENTO DE
INGENIERÍA ELÉCTRICA, ELECTRÓNICA, AUTOMÁTICA Y COMUNICACIONES o DEPARTAMENTO DE MATEMÁTICAS o cualquier otro
de la UCLM al que pertenezca el director. \par }
\par}

\vskip 3em
{\centering\bfseries\large
ANTEPROYECTO DEL TRABAJO FIN DE GRADO
\par}

{\centering\bfseries\large
GRADO EN INGENIERÍA INFORMÁTICA
\par}

{\centering\bfseries\large
TECNOLOGÍA ESPECÍFICA DE / INTENSIFICACIÓN / ITINERARIO DE \footnote{INGENIER\'IA DEL SOFTWARE o INGENIER\'IA DE COMPUTADORES o COMPUTACI\'ON o
TECNOLOG\'IAS DE LA INFORMACI\'ON (esta \'ultima est\'a tambi\'en asociada a los TFG del \textbf{curso} de
\textbf{adaptaci\'on})}
\par}

\vskip 3em
{\centering\bfseries
(Título del TFG)
\par}

\vfill % Rellena espacio automáticamente hasta ajustar al margen inferior 

Autor: Fernando Luján Martínez

Director: Gregorio Díaz Descalzo
\footnote{Sólo en el caso de que haya un segundo director.} (nombre y apellidos)

\vskip 1em 
\mbox{\begin{minipage}[b][2.5cm][c]{0.7\linewidth} {\large Mes, 2018}  \end{minipage}}
\hfill % Rellena espacio automáticamente hasta que la image se ajuste a la derecha
\mbox{\includegraphics[width=2.5cm]{logo-euroinf.png}}
\vskip 1em %Un línea en blanco para que el logo de Euroinf no se pegue al footnote

\clearpage\pagestyle{empty}
\thispagestyle{FirstPage}

\setcounter{tocdepth}{10}
\renewcommand\contentsname{\'Indice de contenido}
\tableofcontents

\clearpage
\pagestyle{empty}

El anteproyecto recoger\'a, en un \textcolor[rgb]{0.5019608,0.0,0.0}{m\'aximo de 10 p\'aginas}, los siguientes
apartados:

\liststyleWWNumxxix
\begin{itemize}
\item Introducci\'on (muy recomendable aunque no obligatorio)
\item Tecnolog\'ia espec\'ifica cursada por el alumno
\item Objetivos
\item M\'etodo y fases de trabajo
\item Medios que se pretenden utilizar
\item Bibliograf\'ia b\'asica consultada en la elaboraci\'on del anteproyecto
\item Contrato de propiedad intelectual (si lo hubiera)
\end{itemize}

\section{1.\ INTRODUCCI\'ON}

El cap\'itulo de introducci\'on podr\'a abordar los siguientes aspectos:

\liststyleWWNumxxix
\begin{itemize}
\item Introducci\'on al tema, entorno en el que el trabajo desempe\~nar\'a su objetivo, \ justificaci\'on de la
importancia del trabajo abordado.
\item Motivaci\'on y antecedentes (con algunas referencias bibliogr\'aficas).
\item Descripci\'on gr\'afica del proyecto (es aconsejable incorporar una figura que describa el trabajo a desarrollar y
que mejore la comprensi\'on del mismo).
\end{itemize}

En la actualidad hay una gran cantidad de 


\section{2. TECNOLOG\'IA ESPEC\'IFICA / INTENSIFICACIÓN / ITINERARIO CURSADO POR EL ALUMNO}
El Trabajo Fin de Grado (TFG, de ahora en adelante) siempre deber\'a demostrar la aplicaci\'on de las competencias
generales de la titulaci\'on. Adem\'as, el TFG deber\'a aplicar \textbf{algunas} de las competencias espec\'ificas
asociadas a la \textbf{Tecnolog\'ia Espec\'ifica o Intensificaci\'on} que el alumno ha cursado. Por lo tanto, el alumno
incluir\'a en el anteproyecto \textbf{dos tablas}. Una tabla para seleccionar la tecnolog\'ia cursada y en la que se
contextualiza el TFG:


\bigskip

{\centering\bfseries
Tabla 1. Tecnolog\'ia Espec\'ifica cursada por el alumno
\par}

\begin{center}
\tablefirsthead{}
\tablehead{}
\tabletail{}
\tablelasttail{}
\begin{supertabular}{m{6.803cm}}
 \rowcolor{gray!50}
{\bfseries Marcar la Tecnolog\'ia Cursada}
\\\hline
~

Tecnolog\'ias de la Informaci\'on\\
\rowcolor{gray!30}
~

Computaci\'on\\
~

Ingenier\'ia del Software\\
\rowcolor{gray!30}
~

Ingenier\'ia de Computadores\\\hline
\end{supertabular}
\end{center}



\clearpage\pagestyle{plain}
\thispagestyle{FirstPage}

\ \ En la segunda tabla, el alumno deber\'a justificar c\'omo \textbf{algunas} de las competencias espec\'ificas de la
intensificaci\'on se aplicar\'an o tomar\'an forma en el TFG, \textbf{La relaci\'on de competencias por
intensificaci\'on se encuentran en el Anexo I al final de este documento. }


\bigskip

{\centering\bfseries
Tabla 2. Justificaci\'on de las competencias espec\'ificas abordadas en el TFG
\par}

\begin{center}
\tablefirsthead{}
\tablehead{}
\tabletail{}
\tablelasttail{}
\begin{supertabular}{m{8.849cm}m{7.605cm}}
\rowcolor{gray!50}
{\color{black} \textbf{Competencias}} &
{\color{black} \textbf{Justificaci\'on}}\\\hline
{\color{black} Competencia 1} &
{\color{black} [Exponer y argumentar c\'omo y en qu\'e parte se va a abordar esta competencia en el TFG]}\\
~
 &
~
\\
\rowcolor{gray!30}
Competencia 2
~
 &
~
[Exponer y argumentar c\'omo y en qu\'e parte se va a abordar esta competencia en el TFG]
\\
~
 &
~
\\\hline
\end{supertabular}
\end{center}

\bigskip

\section{3. \ OBJETIVOS}
De acuerdo a la Introducci\'on, el alumno deber\'a especificar cu\'al o cuáles son las hip\'otesis de trabajo de las
que se parten, qu\'e se pretende resolver, y en base a eso formular el objetivo principal del TFG.

El objetivo principal deber\'a desglosarse en sub-objetivos parciales. Los sub-objetivos deberán describirse de forma breve y concisa.

Como pre\'ambulo a la formulaci\'on del objetivo parcial, el alumno deber\'a discutir sobre las limitaciones y
condicionantes a tener en cuenta en el desarrollo del TFG (lenguaje de desarrollo, equipos, madurez de la tecnolog\'ia,
etc.).

Del mismo modo, ser\'a recomendable incluir una lista preliminar de requisitos del sistema a construir.

\section[4. M\'ETODO Y FASES DE TRABAJO]{4. M\'ETODO Y FASES DE TRABAJO}

\bigskip

Para el desarrollo del proyecto, el alumno deber\'a seguir alg\'un proceso o metodolog\'ia af\'in al problema que
pretende resolver. Para ello, deber\'a aportar una peque\~na descripci\'on del proceso o metodolog\'ia (no m\'as de una
p\'agina) y \textbf{justificar su adecuaci\'on al problema a resolver}.

Del mismo modo, el alumno podrá realizar una breve planificaci\'on de la ejecuci\'on del proyecto seg\'un el proceso
o metodolog\'ia seleccionada.

Como parte de la descripci\'on del m\'etodo y las fases de trabajo, el alumno podr\'a incluir una descripci\'on
preliminar de las tareas, una planificaci\'on temporal, diagramas de Gantt o recursos similares que pueda considerar
necesarios.

Si hubiera m\'as de una metodolog\'ia que a juicio del alumno podr\'ia ser af\'in al proyecto, \'estas deber\'an
mencionarse, y justificar la que considera m\'as adecuada (esto puede considerarse parte de la justificaci\'on a la
adecuaci\'on al problema a resolver).

\section{5. MEDIOS QUE SE PRETENDEN UTILIZAR}
\section{5.1. Medios Hardware}

El alumno deber\'a describir los medios hardware que prev\'e ser\'an necesarios para el desarrollo del proyecto.

\section{5.2. Medios Software}

El alumno deber\'a describir los medios software (lenguajes, entornos de desarrollo, herramientas de gesti\'on y
planificaci\'on, etc.) que prev\'e ser\'an necesarios para el desarrollo del proyecto

%\clearpage\pagestyle{empty}
%\thispagestyle{FirstPage}

\section[6. REFERENCIAS]{6. REFERENCIAS}
En esta secci\'on se incluir\'an todas las referencias bibliogr\'aficas, ordenadas alfab\'eticamente por el primer
apellido del primer autor, de las obras de las cuales se haya realizado alguna cita en los apartados anteriores. Las
referencias deber\'an contener datos b\'asicos como nombre y apellidos de los autores, t\'itulo de la obra, evento al
que pertenece, p\'aginas, fecha y lugar de celebraci\'on (si se tratara de art\'iculos de congreso), ISBN, editorial y
ciudad (si se tratara de libro), nombre de revista, p\'aginas, volumen y n\'umero (si se tratara de revista), etc.

Se emplear\'a un formato de referencia reconocido en el \'ambito acad\'emico como
ACM\footnote{http://www.acm.org/sigs/publications/proceedings-templates}\footnote{http://www.cs.ucy.ac.cy/\~{}chryssis/specs/ACM-refguide.pdf}.
Otros formatos aconsejables son, por ejemplo, IEEE, AMA, APA y AMA.

Ejemplos de referencias con formato ACM:

\liststyleLi
\begin{itemize}
\item Para un art\'iculo de revista:
\end{itemize}
Bowman, M., Debray, S. K., and Peterson, L. L. 1993. Reasoning about naming systems. \textit{ACM Trans. Program. Lang.
Syst.} 15, 5 (Nov. 1993), 795-825. DOI= \url{http://doi.acm.org/10.1145/161468.16147}.

\liststyleLii
\begin{itemize}
\item Para un informe t\'ecnico
\end{itemize}
Ding, W. and Marchionini, G. 1997. \textit{A Study on Video Browsing Strategies. Technical Report}. University of
Maryland at College Park.

\liststyleLiii
\begin{itemize}
\item Para un libro

Tavel, P. 2007. \textit{Modeling and Simulation Design.} AK Peters Ltd., Natick, MA, USA.
\end{itemize}
\liststyleLiv
\begin{itemize}
\item Para un cap\'itulo de libro:
\end{itemize}
Greiner, R. 1999. Explanation-based learning. In Wilson and F. Keil, R. eds. \textit{The Encyclopedia of Cognitive
Science}, MIT Press, Cambridge, MA, USA. 301-303.

\liststyleLv
\begin{itemize}
\item Para un art\'iculo en las actas de un congreso:
\end{itemize}
Fr\"ohlich, B. and Plate, J. 2000. The cubic mouse: a new device for three-dimensional input. In \textit{Proceedings of
the SIGCHI Conference on Human Factors in Computing Systems} (The Hague, The Netherlands, April 01 - 06, 2000). CHI
'00. ACM, New York, NY, 526-531. DOI= \url{http://doi.acm.org/10.1145/332040.332491}.

\liststyleLvi
\begin{itemize}
\item Para un p\'agina de Internet (con autores conocidos)

Steele, B. Look, Ma, no wires! Cornell class project tests wireless networking, \textit{Cornell Chronicle, 31 }(35).
Retrieved February 15, 2004, from Columbia University:
\url{http://www.news.cornell.edu/Chronicle/00/5.18.00/wireless_class.html}
\item Para un p\'agina de Internet (con autores desconocidos)

MIT Project Oxygen: Overview, 2004. Retrieved March 15, 2005, from Computer Science and Artificial Intelligence
Laboratory, Massachusetts Institute of Technology: \url{http://oxygen.lcs.mit.edu/Overview.html}. \
\end{itemize}

\bigskip

\section{7. CONTRATO DE PROPIEDAD INTELECTUAL (si lo hubiera)}

\bigskip

\clearpage\clearpage\pagestyle{plain}
\thispagestyle{FirstPage}
{\centering
\textstyleStrongEmphasis{ANEXO I: Descripci\'on de Competencias por Intensificaci\'on o Tecnolog\'ia
Espec\'ifica}\footnote{Este anexo se deber\'a borrar y no deber\'a ser incluido en el documento de anteproyecto final}
\par}

{\bfseries
Intensificaci\'on de Computaci\'on:}

\liststyleLvii
\begin{itemize}
\item Capacidad para tener un conocimiento profundo de los principios fundamentales y modelos de la computaci\'on y
saberlos aplicar para interpretar, seleccionar, valorar, modelar, y crear nuevos conceptos, teor\'ias, usos y
desarrollos tecnol\'ogicos relacionados con la inform\'atica.
\item Capacidad para conocer los fundamentos te\'oricos de los lenguajes de programaci\'on y las t\'ecnicas de
procesamiento l\'exico, sint\'actico y sem\'antico asociadas, y saber aplicarlas para la creaci\'on, dise\~no y
procesamiento de lenguajes.
\item Capacidad para evaluar la complejidad computacional de un problema, conocer estrategias algor\'itmicas que puedan
conducir a su resoluci\'on y recomendar, desarrollar e implementar aquella que garantice el mejor rendimiento de
acuerdo con los requisitos establecidos.
\item Capacidad para conocer los fundamentos, paradigmas y t\'ecnicas propias de los sistemas inteligentes y analizar,
dise\~nar y construir sistemas, servicios y aplicaciones inform\'aticas que utilicen dichas t\'ecnicas en cualquier
\'ambito de aplicaci\'on.
\item Capacidad para adquirir, obtener, formalizar y representar el conocimiento humano en una forma computable para la
resoluci\'on de problemas mediante un sistema inform\'atico en cualquier \'ambito de aplicaci\'on, particularmente los
relacionados con aspectos de computaci\'on, percepci\'on y actuaci\'on en ambientes entornos inteligentes.
\item Capacidad para desarrollar y evaluar sistemas interactivos y de presentaci\'on de informaci\'on compleja y su
aplicaci\'on a la resoluci\'on de problemas de dise\~no de interacci\'on persona computadora.
\item Capacidad para conocer y desarrollar t\'ecnicas de aprendizaje computacional y dise\~nar e implementar
aplicaciones y sistemas que las utilicen, incluyendo las dedicadas a extracci\'on autom\'atica de informaci\'on y
conocimiento a partir de grandes vol\'umenes de datos.
\end{itemize}
{\bfseries
Intensificaci\'on de Ingenier\'ia de Computadores:}

\liststyleLviii
\begin{itemize}
\item Capacidad de dise\~nar y construir sistemas digitales, incluyendo computadores, sistemas basados en
microprocesador y sistemas de comunicaciones.
\item Capacidad de desarrollar procesadores espec\'ificos y sistemas empotrados, as\'i como desarrollar y optimizar el
software de dichos sistemas.
\item Capacidad de analizar y evaluar arquitecturas de computadores, incluyendo plataformas paralelas y distribuidas,
as\'i como desarrollar y optimizar software para las mismas.
\item Capacidad de dise\~nar e implementar software de sistema y de comunicaciones.
\item Capacidad de analizar, evaluar y seleccionar las plataformas hardware y software m\'as adecuadas para el soporte
de aplicaciones empotradas y de tiempo real.
\item Capacidad para comprender, aplicar y gestionar la garant\'ia y seguridad de los sistemas inform\'aticos.
\item Capacidad para analizar, evaluar, seleccionar y configurar plataformas hardware para el desarrollo y ejecuci\'on
de aplicaciones y servicios inform\'aticos.
\item Capacidad para dise\~nar, desplegar, administrar y gestionar redes de computadores.
\end{itemize}

\clearpage\pagestyle{plain}
\thispagestyle{FirstPage}
{\bfseries
Intensificaci\'on de Ingenier\'ia del Software:}

\liststyleLix
\begin{itemize}
\item Capacidad para desarrollar, mantener y evaluar servicios y sistemas software que satisfagan todos los requisitos
del usuario y se comporten de forma fiable y eficiente, sean asequibles de desarrollar y mantener y cumplan normas de
calidad, aplicando las teor\'ias, principios, m\'etodos y pr\'acticas de la Ingenier\'ia del Software.
\item Capacidad para valorar las necesidades del cliente y especificar los requisitos software para satisfacer estas
necesidades, reconciliando objetivos en conflicto mediante la b\'usqueda de compromisos aceptables dentro de las
limitaciones derivadas del coste, del tiempo, de la existencia de sistemas ya desarrollados y de las propias
organizaciones.
\item Capacidad de dar soluci\'on a problemas de integraci\'on en funci\'on de las estrategias, est\'andares y
tecnolog\'ias disponibles.
\item Capacidad de identificar y analizar problemas y dise\~nar, desarrollar, implementar, verificar y documentar
soluciones software sobre la base de un conocimiento adecuado de las teor\'ias, modelos y t\'ecnicas actuales.
\item Capacidad de identificar, evaluar y gestionar los riesgos potenciales asociados que pudieran presentarse.
\item Capacidad para dise\~nar soluciones apropiadas en uno o m\'as dominios de aplicaci\'on utilizando m\'etodos de la
ingenier\'ia del software que integren aspectos \'eticos, sociales, legales y econ\'omicos.
\end{itemize}
{\bfseries
Intensificaci\'on de Tecnolog\'ias de la Informaci\'on:}

\liststyleLx
\begin{itemize}
\item Capacidad para comprender el entorno de una organizaci\'on y sus necesidades en el \'ambito de las tecnolog\'ias
de la informaci\'on y las comunicaciones.
\item Capacidad para seleccionar, dise\~nar, desplegar, integrar, evaluar, construir, gestionar, explotar y mantener las
tecnolog\'ias de hardware, software y redes, dentro de los par\'ametros de coste y calidad adecuados.
\item Capacidad para emplear metodolog\'ias centradas en el usuario y la organizaci\'on para el desarrollo, evaluaci\'on
y gesti\'on de aplicaciones y sistemas basados en tecnolog\'ias de la informaci\'on que aseguren la accesibilidad,
ergonom\'ia y usabilidad de los sistemas.
\item Capacidad para seleccionar, dise\~nar, desplegar, integrar y gestionar redes e infraestructuras de comunicaciones
en una organizaci\'on.
\item Capacidad para seleccionar, desplegar, integrar y gestionar sistemas de informaci\'on que satisfagan las
necesidades de la organizaci\'on, con los criterios de coste y calidad identificados.
\item Capacidad de concebir sistemas, aplicaciones y servicios basados en tecnolog\'ias de red, incluyendo Internet,
web, comercio electr\'onico, multimedia, servicios interactivos y computaci\'on m\'ovil.
\item Capacidad para comprender, aplicar y gestionar la garant\'ia y seguridad de los sistemas inform\'aticos.
\end{itemize}

\bigskip
\end{document}
